\section*{Higham Chapter 6, Problem 6.8: Lean Triangularization Bridge}

I am formalizing Higham's finite-dimensional complex matrix exercise:

\[
\forall A\in\mathbb C^{n\times n},\ \forall \delta>0,\ \exists\hbox{ a
consistent/subordinate matrix norm } \|\cdot\| \hbox{ such that }
\|A\| \le \rho(A)+\delta .
\]

The analytic weighted-infinity part is already checked in Lean once an upper
triangular coordinate model is supplied.  The existing checked theorem consumes
the following data:

\[
T = P^{-1} A P,\qquad T\hbox{ upper triangular},\qquad
\forall i,\ |T_{ii}|\le \rho ,
\]

and returns a subordinate matrix norm bounded by \(\rho+\delta\).  The remaining
blocker is the algebraic bridge from an arbitrary complex matrix to such a
triangular model, plus the diagonal-to-spectral-radius bound.

Current mathlib facts I found locally:

\begin{itemize}
\item \texttt{Mathlib.LinearAlgebra.Eigenspace.Triangularizable} explicitly
has TODO text saying triangularizable endomorphisms are not yet defined as a
maximal invariant flag/basis theorem.
\item It does have \texttt{Module.End.exists\_eigenvalue}: every nontrivial
finite-dimensional endomorphism over an algebraically closed field has an
eigenvalue.
\item It has \texttt{Module.End.iSup\_maxGenEigenspace\_eq\_top} and
\texttt{Module.End.genEigenspace\_restrict\_eq\_top}.
\item \texttt{Mathlib.LinearAlgebra.Basis.Flag} has \texttt{Basis.flag},
\texttt{flag\_le\_iff}, \texttt{flag\_succ}, \texttt{self\_mem\_flag},
and \texttt{flag\_le\_ker\_coord}.
\item Matrices use \texttt{Matrix.BlockTriangular id} for upper triangularity;
\texttt{Matrix.charpoly\_of\_upperTriangular} exists.
\item I did not find a ready-made declaration producing a basis whose
\texttt{LinearMap.toMatrix b b f} is upper triangular.
\end{itemize}

Please give a Lean-friendly implementation plan for the smallest bridge:

\[
\exists b : \mathrm{Basis}\ (\mathrm{Fin}\ n)\ \mathbb C\ (\mathrm{Fin}\ n\to\mathbb C),
\quad \mathrm{Matrix.BlockTriangular}\ \mathrm{id}\
  (\mathrm{LinearMap.toMatrix}\ b\ b\ f),
\]

or an equivalent local predicate that is easier to prove and enough for the
weighted-norm theorem.

Specific questions:

\begin{enumerate}
\item If current mathlib has a hidden ready-made basis/flag theorem I missed,
name the exact declaration(s).
\item If not, propose the smallest induction proof over \(n\).  I want exact
submodule/basis-extension lemmas to search for: how to choose an eigenvector,
extend it to a basis, pass to the quotient or invariant complement/tail, and
prove the resulting \texttt{toMatrix} is upper triangular.
\item Is it easier to use generalized eigenspaces from
\texttt{iSup\_maxGenEigenspace\_eq\_top}, or ordinary eigenvector induction?
\item How should the diagonal entry be connected to an eigenvalue modulus for
my local predicate
\[
\mathrm{ComplexVectorMapEigenvalueModulusSet}(T)
 = \{\,|\lambda| : \exists x\ne0,\ T x=\lambda x\,\},
\]
so that a supplied maximum \(\rho\) gives \(|T_{ii}|\le\rho\)?
\item Flag zero-dimensional and nonempty-dimensional caveats, and whether it is
better to state the bridge as a new local structure containing the basis,
upper-triangular matrix, diagonal eigenvalue witnesses, and similarity action
rather than as a theorem immediately producing \(P^{-1}AP\).
\end{enumerate}

Please answer with concrete Lean theorem shapes and proof-route details.  Do
not rely on unproved Schur/unitary triangularization; ordinary algebraic
triangularization is enough.
